% ============================================================
% Visualization Guide for Project 3 (PDE IBVP)
% Compile: tectonic -X compile visualization_guide_ibvp.tex
% ============================================================
\documentclass[11pt]{article}
\usepackage[a4paper,margin=2.3cm]{geometry}
\usepackage{amsmath,amssymb}
\usepackage{booktabs}
\usepackage{fontspec}
\usepackage{microtype}
\usepackage{xcolor}
\usepackage{enumitem}
\usepackage{hyperref}
\setmainfont{Times New Roman}
\newcommand{\e}[1]{{\emojifont #1}}
\setlength{\parindent}{0pt}
\setlength{\parskip}{0.5em}
\renewcommand{\arraystretch}{1.3}
\begin{document}
\begin{center}
  {\LARGE\bfseries Visualization Guide for Project 3 (PDE IBVP)}\\[0.4em]
  {\large Companion to the heat and hyperbolic scaffolds}
\end{center}
\vspace{0.2em}
\hrule
\vspace{0.6em}

The IBVP project produces four classes of figures: \textbf{time-series}
(convergence + amplitude + energy drift), \textbf{heatmaps / colormaps}
(space--time profiles, especially for diffusion and transport),
\textbf{animations} (long-time wave / shock evolution), and the
\textbf{cost--accuracy} plot that carries the module's efficiency outcome.
This guide
follows the same Rule 1--6 conventions as the ODE/BVP guides, with
additions for time-evolution and shock-safe colour maps; Rule 7 is
specific to Project 3, because this is the project on which learning
outcome D is assessed.

% ----------------------------------------------------------------------
\section*{Rule 1 --- One claim per figure, with a captioned interpretation}

Every figure has a \textbf{title} (the conclusion, not the content) and
a \textbf{caption} that says what is shown \emph{and} why it matters.

\begin{itemize}[leftmargin=2em, itemsep=1pt]
  \item\textbf{Time-series (parabolic):}
        title like \emph{``$L^2$ error decays at order 1.5 in time''};
        axes $h$ (or $\Delta t$) and $L^2$ error; legend distinguishes schemes.
  \item\textbf{Energy / amplitude (wave):}
        title like \emph{``leapfrog conserves the staggered discrete energy
        to roundoff over 100 periods (drift $3\times10^{-13}$)''};
        axes time and $E^{n+1/2}/E^{1/2}$. Roundoff-scale conservation is
        an implementation check, \emph{not} an accuracy statement: the gap
        between the discrete and the continuous energy
        $E=\tfrac12\int(u_t^2+c^2u_x^2)\,dx$ decreases at the scheme's
        order in $h$, and is reported as a separate number.
  \item\textbf{Heatmap (space--time):}
        title like \emph{``heat smooths the kink''} (mass is \emph{not}
        conserved under zero Dirichlet data --- the default sine solution
        loses mass as it decays; claim conservation only with periodic or
        no-flux boundaries, and include boundary fluxes in the mass-balance
        diagnostic for open systems);
        axes $x$ and $t$, colourmap \texttt{viridis} (NOT \texttt{jet}).
  \item\textbf{Dispersion curve:}
        title like \emph{``leapfrog phase speed lags at moderate $kh$''};
        axes $kh$ and $\omega_h(k)/ck$, reference line $y=1$.
\end{itemize}

% ----------------------------------------------------------------------
\section*{Rule 2 --- Colour maps that do not hide structure}

\begin{itemize}[leftmargin=2em, itemsep=1pt]
  \item For \textbf{heatmaps} use \texttt{viridis}, \texttt{plasma}, or
        \texttt{coolwarm} (perceptually uniform; colour-blind safe).
  \item For \textbf{diverging fields} (residual, signed error) use
        \texttt{coolwarm} centred on zero; do not use \texttt{jet} (it
        exaggerates differences and fails in print).
  \item For \textbf{time-series} prefer \textbf{colour-blind-safe pairs}:
        \texttt{tab:blue}/\texttt{tab:orange} or solid + dashed.
\end{itemize}

% ----------------------------------------------------------------------
\section*{Rule 3 --- Reproduction card: the exact script + commit}

Every reported number has a pinned commit and a one-line run command.
\textit{Done} $=$ the command prints a line that matches the report
within the stated \emph{diagnostic-specific} tolerance --- e.g.\ the
fitted order within $\pm0.05$ of the reported slope, a conservation
drift below the declared roundoff threshold, a PASS/FAIL verdict
against the tolerance stated in the report --- not a universal digit
count. Each skeleton is itself a runnable driver: after its TASKs are
implemented, \texttt{python heat\_equation\_skeleton.py} prints the heat
centre value with a PASS/FAIL verdict, and
\texttt{python hyperbolic\_skeleton.py} prints the wave centre value with a
PASS/FAIL verdict. The drivers check against the analytic solutions inside
the files, so a reproduced number needs no instructor solution.

% ----------------------------------------------------------------------
\section*{Rule 4 --- Axes, units, legends, reproducibility}

\begin{itemize}[leftmargin=2em, itemsep=1pt]
  \item Every axis carries a unit: $h$ [m], $\Delta t$ [s], $x$ [m], $t$ [s], $T$ [K], ...
  \item Every curve has a legend entry; lines that overlap are differentiated
        by linestyle (solid / dashed / dotted) and marker.
  \item Random seeds are fixed and recorded; the figure caption
        states the seed.
  \item Generated figures follow the module's single commit policy:
        \textbf{final} figures that appear in the report (and the script
        that produced them) are committed; intermediate or disposable
        outputs (scratch plots, caches, old \texttt{.png} from earlier
        runs) are not. Keep final figures in \texttt{figures/}, source in
        \texttt{code/}; rebuild before committing so no stale output is
        committed.
\end{itemize}

% ----------------------------------------------------------------------
\section*{Rule 5 --- Fidelity diagnostics are the figure's spine}

A fidelity figure is not \emph{``a picture of the solution''} but
\emph{``a measurement of a property the continuous PDE has''} that a
wrong method would break. Pick the diagnostic first, then design the run
that produces it.

\begin{itemize}[leftmargin=2em, itemsep=1pt]
  \item\textbf{Parabolic:} maximum principle ($\min U \ge 0$, $\max U \le
        \max U_0$) and a non-increasing discrete energy --- \emph{under the
        data assumptions that guarantee them}, which you state explicitly
        (e.g.\ non-negative initial data, no interior source, and boundary
        data within the initial range; a dissipative boundary condition for
        the energy); check them only on runs where those assumptions hold
        --- plus a \emph{time-discrete} residual of the evolution
        equation, e.g.\ for backward Euler
        \[
        \Big\|\frac{U^{n+1}-U^n}{\Delta t} - L_h U^{n+1} - f^{n+1}\Big\|_\infty,
        \]
        \emph{not} the elliptic residual $\|L_h U - f\|_\infty$, which can be
        nonzero for an exact transient solution and would falsely flag a
        solver defect. Test discretisation accuracy with an independent
        reference or error diagnostic instead.
  \item\textbf{Wave:} relative drift of the staggered discrete energy
        $|E^{n+1/2} - E^{1/2}|/E^{1/2}$ over 100 periods (roundoff scale:
        an implementation check, not accuracy), the discrete-to-continuous
        gap $|E^{1/2} - E_{\mathrm{cont}}|$ decreasing at the scheme's
        order in $h$ (for $u=\sin\pi x\cos\pi t$, $c=1$:
        $E_{\mathrm{cont}}=\pi^2/4$), phase displacement per wavelength,
        amplitude decay per period.
  \item\textbf{Transport / conservation law:} weighted total mass
        $M_h=\sum_j w_j U_j$, with the weights matched to what $U_j$
        stores: cell widths or volumes when $U_j$ is a finite-volume cell
        average, trapezoidal weights (half at the ends) when $U_j$ is a
        nodal value on a suitable grid, and its change compared against
        boundary fluxes and sources ---
        $|M_h(t)-M_h(0)|$ alone rises legitimately on an open or forced
        problem and would mislabel conservation; with periodic or no-flux
        boundaries and no source it is the balance residual. Also: shock
        position vs.\ analytic, TVD property (no new extrema), entropy
        residual.
\end{itemize}

% ----------------------------------------------------------------------
\section*{Rule 6 --- Time-evolution specifics for IBVP}

The IBVP project stresses time-evolution. The list below is
\emph{pathway-dependent}: apply the items that fit your PDE and scheme,
and state which apply. A diffusion project has no dispersion curve, and
a scheme whose \emph{unconditional stability for this operator has been
established} (state the method, the discrete operator, and the
assumptions, e.g.\ backward Euler or Crank--Nicolson for the heat
equation $u_t = u_{xx}$ with $L_h$ the symmetric negative-semidefinite
discretisation of $\partial_{xx}$ appearing in the Rule 5 residual) has
no stability boundary to sweep; a wave or transport project must not
omit its energy/mass and phase diagnostics.

\begin{itemize}[leftmargin=2em, itemsep=1pt]
  \item\textbf{Log--log refinement plot} for the observed order:
        $L^2$ error vs $h$ (or $\Delta t$) on log--log axes; fit a
        straight line, record the slope $\pm$ standard error.
        \emph{(All pathways.)}
  \item\textbf{Energy (or mass) drift plot} for waves / transport:
        $E^{n+1/2}/E^{1/2}$ or $M_h(t)/M_h(0)$ vs $t$ on a \emph{linear}
        $y$-axis (so a 1\% drift is visible) and, when the ratio stays
        positive, a \emph{log} $y$ panel (so exponential behaviour is
        visible). If the reference value vanishes ($E^{1/2}=0$ or
        $M_h(0)=0$) the ratio is undefined --- plot the absolute drift
        $|E^{n+1/2}-E^{1/2}|$ or $M_h(t)$ itself instead. Two panels, one
        figure. \emph{(Wave and transport pathways.)}
  \item\textbf{Dispersion relation} $\omega_h(k)$ vs analytic
        $\omega(k)=ck$ for at least 4 $k$ values spanning the resolvable
        range; the offset per wavelength is the report's number.
        \emph{(Wave pathway --- an advection or heat project may instead
        report a phase-displacement or decay measurement; derive what
        applies.)}
  \item\textbf{Stability-boundary sweep} for the topic's analysis:
        run just below and just above the predicted bound, record the
        measured amplification at the most-dangerous mode.
        \emph{(Only for a scheme whose stability is conditional --- a
        derived bound exists to sweep, e.g.\ explicit FTCS or leapfrog at
        the relevant $r$ or CFL value. A scheme whose unconditional
        stability for this operator has been established, stated with its
        assumptions, has no boundary to sweep.)}
\end{itemize}

% ----------------------------------------------------------------------
\section*{Rule 7 --- Cost--accuracy evidence (Learning Outcome D)}

Project 3 is where the module's efficiency outcome (LO D: reflect on
computational efficiency in terms of cost and accuracy) is assessed, so one
figure or table must carry it. \emph{Accuracy alone is not the claim; the
claim is accuracy bought per unit of work.}

\begin{itemize}[leftmargin=2em, itemsep=1pt]
  \item\textbf{Cost--accuracy plot:} $x$-axis = cost, $y$-axis = achieved
        error, one curve per scheme, both reaching the \emph{same} accuracy
        target. A second panel with a log $y$-axis makes the asymptotic
        separation visible; a single linear panel is acceptable when only
        two operating points are compared.
  \item\textbf{Name the cost measure} and use it consistently: wall-clock
        time [s] (per \texttt{time.perf\_counter} around the solve only,
        \emph{excluding} figure generation and I/O), iteration count for an
        iterative solve, or operation count (e.g.\ one matvec $=O(N)$ per
        step). Wall-clock is machine-dependent --- state the machine and
        repeat timing runs, or report iteration/operation counts alongside it.
  \item\textbf{Match the accuracy target} before comparing cost: an explicit
        scheme run at a smaller $\Delta t$ ``wins'' on error only because it
        was given more work. Either fix the target error and report the cost
        each scheme needs to reach it, or plot the whole error-vs-cost curve
        and read the comparison off the curves.
  \item\textbf{State what you did \emph{not} count} (setup, factorisation,
        precomputation, figure rendering) and hold it constant across schemes.
  \item\textbf{The honest negative result is still a result.} ``Scheme A
        reaches $10^{-4}$ in 0.8\,s while scheme B needs 3.1\,s, so A is
        cheaper \emph{for this problem and this target}'' is a complete LO D
        answer; a claim of ``A is more efficient'' with no measured cost is not.
\end{itemize}

% ----------------------------------------------------------------------
\section*{Time-series vs heatmap vs animation: which to use when}

\begin{tabular}{@{}p{0.25\linewidth}p{0.30\linewidth}p{0.38\linewidth}@{}}
\toprule
\textbf{Story you want to tell} & \textbf{Best figure} & \textbf{Why} \\
\midrule
``The error decays at order $p$'' & log--log refinement & slope $=$ order; cited in the report \\
``The method damps / oscillates / preserves the wave'' & $E(t)/E(0)$ vs $t$ & visible on a linear $y$ \\
``The shock is at the right place and right height'' & $u(x,T)$ profile (line + overlay) & compares to analytic Riemann \\
``Long-time behaviour is stable'' & animation (MP4 / GIF) + drift plot & animation shows motion; the companion drift plot turns ``stable'' into a number \\
``This scheme buys accuracy more cheaply'' & error vs.\ cost at a matched accuracy target & turns LO D into a measured number instead of an assertion \\
``Stability boundary at $r=1/2$ holds / breaks'' & blow-up montage + amplification & one figure, one claim \\
\bottomrule
\end{tabular}

% ----------------------------------------------------------------------
\section*{The ``what is wrong with this figure'' checklist}

Before you commit a figure, scan it for these common failures:
\begin{itemize}[leftmargin=2em, itemsep=1pt]
  \item Rainbow / jet colormap on a heatmap $\to$ \textbf{change to
        \texttt{viridis}} (perceptually uniform, print-safe).
  \item Smooth sine that decays through zero on a linear $y$-axis $\to$
        its signed values cannot go on a $\log$ axis (which has no zero
        and discards negative data); plot the positive envelope
        $\max_x|u(x,t)|$ (or the $\|u(\cdot,t)\|_2$ norm) on $\log y$ to
        expose the decay, and keep the signed trace on the linear axis.
  \item Title says ``solution'' or ``result'' $\to$ \textbf{rewrite} the
        title as the conclusion, e.g.\ \emph{``$L^2$ error is
        $O(h^{1.97})$ for CN''}.
  \item Caption only describes what's on the figure $\to$ \textbf{add
        the why-it-matters} line: \emph{``the slope matches the second
        order of Crank--Nicolson in time''}.
  \item Animation has no time scale / axis label $\to$ add a frame
        counter and the time stamp; the reader must be able to read
        the evolution.
  \item Stale \texttt{.png} from an earlier commit $\to$ rebuild before
        committing; \texttt{run\_all.py} regenerates everything.
\end{itemize}

\end{document}